
\(\Grb(X)\)
devient ainsi un module sur l'anneau \(\Gro(X)\) (et même un module gradué sur l'anneau gradué \(\Gro(X)\), à condition de changer de signe de la graduation de \(\Grb(X)\) en posant \(\Gr^i_\bullet(X)=\Gr_{-i}(X)\) \((i\in\mathbb Z)\)).
Considérons maintenant un morphisme propre
\[
 f:X'\longrightarrow X
\]
de schémas noethériens. Alors \(f\) induit un homomorphisme
\[
 f^*:\Gro(X)\longrightarrow\Gro(X')
\]
d'anneaux gradués (V), et d'après 1.1.8 un homomorphisme
\[
 f_*:\Grb(X')\longrightarrow\Grb(X)
\]
de groupes abéliens gradués. Alors la formule de projection (IV 2.11) et (3.1.1) donnent par passage aux gradués associés la formule de projection graduée
\begin{equation}\tag{3.1.2}
 x\cdot f_*(y)=f_*\bigl(f^*(x)\cdot y\bigr)
 \qquad (x\in\Gro(X),\ y\in\Grb(X')).
\end{equation}

\subsection*{3.2.}
Soit \(f:X'\to X\) un morphisme projectif d'intersection complète (Exp. VIII) et de dimension virtuelle \(r\) de schémas noethériens, \(X\) ayant un \(\underline{O}_X\)-module inversible ample. Dans Exp.VIII on a démontré que \(f\) induit un homomorphisme
\begin{equation}\tag{3.2.1}
 f_*:\Filt^k(X')_{\mathbb Q}\longrightarrow\Filt^{k-r}(X)_{\mathbb Q}\qquad(k\in\mathbb Z)
\end{equation}
qui satisfait d'ailleurs à la formule de projection (Exp. IV) :
\begin{equation}\tag{3.2.2}
 f_*((f^*(x)\cdot y)=x\cdot f_*(Y)
 \qquad (x\in K^{\bullet}(X)_{\mathbb Q},\ y\in K^{\bullet}(X')_{\mathbb Q}).
\end{equation}
